\hspace*{\fill}
\section{系统需求分析与总体设计}

本章首先明确本系统所需解决的问题与所需交付的功能，进而给出系统总体架构图，并对各模块的职责及其之间的数据流加以说明。

\subsection{功能需求分析}

从教务人员的视角出发，本系统的核心功能涵盖数据维护、自动排课、冲突检测、结果展示与过程可视化五个方面。数据维护要求系统能够查看与维护教师、教室、班级、课程等基础信息，并能清晰呈现它们之间的关联关系（如某门课程由哪位教师讲授、由哪些班级合班授课等）。自动排课要求系统在用户给定的输入数据与算法参数下，自动生成一份满足全部硬约束的课表。冲突检测要求系统对任意候选课表（无论由 GA 生成还是人工导入）逐项指出哪些资源在哪个时间槽发生了冲突、哪节课的教室容量或类型不符合要求。结果展示要求系统能够从班级、教师、教室三个视角查看排课结果，并提供日课时分布热力图、教室占用热力图等概览信息。过程可视化则要求用户在运行 GA 的同时可实时观察适应度的下降曲线、硬冲突数量的变化等指标，以便评估算法的实际收敛情况。

\subsection{非功能需求分析}

非功能需求主要涉及响应速度、易用性、可扩展性与可解释性四个方面。响应速度上，对于一个 12 班 / 39 课的中等规模实例，排课算法应在普通笔记本电脑上数秒内完成。易用性上，界面应保持简洁，教务人员无需 Python 基础即可完成全部交互。可扩展性上，当输入数据发生变化时，核心算法代码无需改动，仅需替换数据集即可完成迁移。可解释性上，算法不应表现为黑盒，必须能够向用户具体指明哪些课程发生了冲突，便于后续排查与处理。

\subsection{系统总体架构设计}

根据上述需求，本系统被划分为六层结构，如图~\ref{fig:arch} 所示。最上层为基于 Streamlit 实现的 Web UI；其下分别为 GA 调度引擎与冲突检测器；再下层为适应度评估与编码及算子两个支持模块；最底层为数据模型层。除 UI 之外，其余各模块均不依赖任何 Web 框架，可作为独立的 Python 库供命令行脚本调用。

\begin{figure}[H]
  \centering
  \begin{tikzpicture}[
    layer/.style={rectangle,draw,thick,rounded corners=3pt,align=center,
                   minimum height=1.1cm,font=\small,inner sep=4pt},
    arrow/.style={-{Latex[length=2mm]},thick,gray!70}
  ]
    \node[layer,fill=blue!10,minimum width=12cm] (ui)
      {\textbf{Streamlit Web UI} \\ \footnotesize 数据展示 \quad 参数配置 \quad 结果渲染};

    \node[layer,fill=green!10,minimum width=5.85cm,below=12pt of ui.south west,anchor=north west] (ga)
      {\textbf{GA 调度引擎} \\ \footnotesize \texttt{scheduler.ga}};
    \node[layer,fill=green!10,minimum width=5.85cm,below=12pt of ui.south east,anchor=north east] (cd)
      {\textbf{冲突检测器} \\ \footnotesize \texttt{scheduler.conflict}};

    \node[layer,fill=orange!12,minimum width=5.85cm,below=12pt of ga.south west,anchor=north west] (ft)
      {\textbf{适应度评估} \\ \footnotesize \texttt{scheduler.fitness}};
    \node[layer,fill=orange!12,minimum width=5.85cm,below=12pt of cd.south east,anchor=north east] (op)
      {\textbf{编码与算子} \\ \footnotesize \texttt{scheduler.operators}};

    \node[layer,fill=red!10,minimum width=12cm,below=12pt of ft.south west,anchor=north west] (dm)
      {\textbf{数据模型层} \\ \footnotesize \texttt{scheduler.models}};

    \draw[arrow] ($(ui.south)+(-3,0)$) -- ($(ga.north)$);
    \draw[arrow] ($(ui.south)+(3,0)$) -- ($(cd.north)$);
    \draw[arrow] (ga.south) -- (ft.north);
    \draw[arrow] (cd.south) -- (op.north);
    \draw[arrow] (ft.south) -- ($(dm.north)+(-3,0)$);
    \draw[arrow] (op.south) -- ($(dm.north)+(3,0)$);
  \end{tikzpicture}
  \caption{系统总体架构}
  \label{fig:arch}
\end{figure}

各模块的职责说明如下。Streamlit Web UI 负责所有人机交互，包括数据展示、参数配置与结果渲染，其内部不持有任何算法逻辑，所有计算均通过调用下层模块完成。GA 调度引擎 (\texttt{scheduler.ga}) 实现遗传算法的主循环，包括种群初始化、代际迭代与精英保留等，对外仅暴露 \texttt{run()} 方法，输入为 \texttt{ProblemInstance} 与 \texttt{GAConfig}，输出为 \texttt{GAResult}。冲突检测器 (\texttt{scheduler.conflict}) 给定一份排课方案后，遍历全部 (时间槽, 资源) 二元组，逐条统计冲突情况，该模块既被适应度函数调用，也被 UI 的"冲突检测报告"页面调用。适应度评估 (\texttt{scheduler.fitness}) 将硬约束违反计数与软约束违反量按权重累加得到 cost，再经 $1/(1+\text{cost})$ 的非线性映射得到适应度 (fitness)。编码与算子 (\texttt{scheduler.operators}) 包含染色体的随机初始化、锦标赛选择、均匀交叉与约束感知变异等算子，均以 \texttt{ProblemInstance} 作为上下文。数据模型层 (\texttt{scheduler.models}) 使用 Python dataclass 对教师、教室、时间槽、班级、课程、课程实例与完整排课方案进行类型化定义，是其余模块所共同依赖的基础。

\subsection{数据建模}

数据模型层所涉及的核心类如表~\ref{tab:models} 所示。所有模型均使用 \texttt{@dataclass(frozen=True)} 装饰，使其在 GA 进化过程中能够被安全共享，避免被无意修改。

\begin{table}[H]
  \centering
  \caption{数据模型类清单}
  \label{tab:models}
  \begin{tabularx}{\linewidth}{l X}
    \toprule
    \textbf{类名} & \textbf{含义与主要字段} \\
    \midrule
    \texttt{Teacher} & 教师：工号、姓名、所在部门、偏好时间段 \\
    \texttt{Classroom} & 教室：编号、名称、容量、类型 (lecture/lab/multimedia) \\
    \texttt{TimeSlot} & 时间槽：id、星期、节次、标签 \\
    \texttt{ClassGroup} & 班级：编号、名称、人数、年级 \\
    \texttt{Course} & 课程：课程号、课程名、任课教师、授课班级、周课时、所需教室类型、最小容量 \\
    \texttt{Lesson} & 课程实例：由 Course 按周课时展开后得到，是 GA 排课的基本单元 \\
    \texttt{ProblemInstance} & 问题实例：聚合上述六类对象，并提供查询方法 \\
    \texttt{Schedule} & 排课方案：按 Lesson 索引存储的 (slot\_idx, room\_idx) 序列 \\
    \bottomrule
  \end{tabularx}
\end{table}

\subsection{典型数据规模}

为给算法明确具体的求解目标，本课题构造了一个相对真实的中等规模实例，其规模如表~\ref{tab:size} 所示，可视为某普通学院一个学期的典型排课情况。该实例的全部数据由 \texttt{scheduler.data.build\_sample\_instance()} 生成，便于复现。

\begin{table}[H]
  \centering
  \caption{测试用例的规模}
  \label{tab:size}
  \begin{tabularx}{\linewidth}{>{\centering\arraybackslash}X >{\centering\arraybackslash}X >{\centering\arraybackslash}X >{\centering\arraybackslash}X >{\centering\arraybackslash}X >{\centering\arraybackslash}X}
    \toprule
    教师数 & 教室数 & 时间槽数 & 班级数 & 课程数 & 课程实例数 \\
    \midrule
    15 & 10 & 50 (5 天 $\times$ 10 节) & 12 & 39 & 62 \\
    \bottomrule
  \end{tabularx}
\end{table}

由上可见，课程实例总数为 62，搜索空间约为 $50^{62} \times 10^{62} \approx 10^{167}$ 量级，完全无法通过枚举求解，必须借助启发式方法。
